% !TEX root = ../memoria.tex

\chapter{Anexo A: Manual de Usuario}
\label{anx:manual_usuario}

Este anexo describe cómo utilizar la aplicación web del gemelo digital una vez
desplegada en local (véase Anexo~\ref{anx:manual_instalacion} para la instalación).
Se asume que el backend está escuchando en \texttt{http://localhost:8080} y el
frontend en \texttt{http://localhost:4200}.

\section{Pantalla principal}
\label{sec:manual_pantalla}

Al abrir la aplicación se carga automáticamente la primera simulación disponible
en el directorio configurado del backend. La pantalla se compone de tres áreas:

\begin{enumerate}
\item \textbf{Cabecera}: muestra el título del proyecto y el botón de cambio de
  tema claro/oscuro de la interfaz.
\item \textbf{Barra lateral izquierda} (\textit{sidebar}): contiene los controles
  principales --- selector de simulación, botón de comparación, parámetros y
  filtro RSRP.
\item \textbf{Mapa principal}: ocupa el resto de la ventana y representa el
  \textit{heatmap} de cobertura sobre la cartografía elegida.
\end{enumerate}

\section{Selección de la simulación activa}
\label{sec:manual_simulacion}

En la sección ``Simulación Activa'' del \textit{sidebar} hay un menú desplegable
con todos los escenarios disponibles. Al elegir uno, el frontend llama al
\textit{endpoint} \texttt{GET /api/simulations/\{id\}}, recibe los puntos de
cobertura y reconstruye el \textit{heatmap}. El mapa se centra automáticamente
en la antena transmisora (TX) del nuevo escenario.

\section{Creación y apoyos de planificación}

Para estudiar un emplazamiento se coloca una antena en el mapa y se introducen
frecuencia, potencia, alturas, área y espaciado. La interfaz solicita el cálculo
al backend Java y presenta la simulación guardada. El resultado representa un
proxy de potencia, no una medición de SS-RSRP.

Como apoyo, las observaciones compatibles se pueden importar para ajustar un
desplazamiento y evaluar sus residuos sobre un conjunto de reserva. La recomendación
evalúa candidatos de un área y propone hasta el número de emplazamientos solicitado
mediante selección voraz;
no garantiza un óptimo global. El patrón direccional, la pérdida manual de
\textit{clutter} y el perfil con Fresnel complementan la exploración. La integración
completa de las interacciones entre estos controles debe verificarse.

\section{Filtro RSRP}
\label{sec:manual_rsrp}

El \textit{slider} ``Umbral RSRP mínimo'' filtra en tiempo real los puntos de
cobertura que se muestran en el \textit{heatmap}. Su rango va de $-140$~dBm
(muestra todos los puntos) a $-40$~dBm. Los nombres de los niveles son etiquetas
visuales del prototipo; el filtro no equivale a un umbral normativo ni acredita
la viabilidad de un servicio. El
contador situado bajo el \textit{slider} indica cuántos puntos pasan el filtro
de un total y permite explorar el mapa con distintos cortes de RSRP.

\section{Inspección puntual}
\label{sec:manual_inspeccion}

Al hacer clic sobre el mapa, el frontend busca el punto de cobertura más cercano
al lugar pulsado. Si la distancia es inferior a aproximadamente 300~m, el
\texttt{InfoPanelComponent} se despliega mostrando: el valor de RSRP en dBm, una
escala visual de barras (1 a 4 segmentos según el intervalo visual del prototipo),
las coordenadas geográficas con seis decimales y los intervalos visuales propios
del prototipo.

\section{Comparación de simulaciones}
\label{sec:manual_comparacion}

El botón ``Comparar simulaciones'' del \textit{sidebar} divide la pantalla en
dos paneles --- A (a la izquierda) y B (a la derecha) --- separados por un
\textit{dashboard} central con métricas. La selección de la simulación A es la
que estuviese activa antes de pulsar el botón; la simulación B se elige en el
desplegable que aparece bajo el botón.

El \textit{dashboard} central muestra:
\begin{itemize}
\item Diferencia media $\overline{\Delta} = \operatorname{media}(\text{RSRP}_A-\text{RSRP}_B)$ sobre los puntos emparejados.
\item Cuando los dos escenarios solapan espacialmente, un desglose de zonas
  donde A es mejor, B es mejor o son similares ($|\Delta| \leq 2$~dB).
\item Distribución comparada de los puntos por intervalos visuales de RSRP.
\item Tabla de parámetros confrontados (frecuencia, potencia, altura, modelo).
\end{itemize}

Cada uno de los dos mapas dispone de su propio selector de cartografía, lo que
permite contrastar la misma cobertura con fondos distintos. El botón
``Cerrar'' devuelve a la vista única.

El emparejamiento espacial usa coordenadas redondeadas a cuatro decimales. Si hay puntos
coincidentes, la diferencia media y el recuento de zonas se calculan sobre esos pares. Si
no hay solapamiento, la interfaz muestra únicamente la diferencia entre las medias globales
de A y B; no presenta recuentos de zonas como si fueran comparaciones punto a punto.

\section{Selector de cartografía}
\label{sec:manual_cartografia}

La tarjeta ``Tema'' en la esquina inferior derecha del mapa expone tres
cartografías base: estándar de OpenStreetMap, satélite de ESRI World Imagery y
relieve de OpenTopoMap. La elección persiste entre sesiones gracias al
almacenamiento local del navegador.

\section{Interpretación del \textit{heatmap}}
\label{sec:manual_interpretacion}

El gradiente cromático del \textit{heatmap} representa los intervalos visuales de
RSRP elegidos por el prototipo, descritos en el Cuadro~\ref{tab:rsrp_umbrales}
(capítulo~\ref{cap:capitulo4}). Sus etiquetas no son categorías universales de
calidad ni predicen por sí solas la disponibilidad de un servicio. El icono de
antena marca la posición exacta del transmisor.

\chapter{Anexo B: Manual de Instalación}
\label{anx:manual_instalacion}

\section{Prerrequisitos}

El sistema se ha desarrollado y probado en Windows~11 con las versiones recogidas
en el Cuadro~\ref{tab:prereqs} del capítulo~\ref{cap:capitulo4}. A modo de
recordatorio, los componentes mínimos necesarios son:

\begin{itemize}
\item \textbf{MATLAB R2025b} (probado con Update~4), solo para reproducir los
  escenarios y referencias MATLAB, con los \textit{toolboxes}
  \textit{Communications}, \textit{Mapping} y \textit{Antenna}. Sin
  \textit{Communications Toolbox} el motor cae en COST-231 con advertencia.
\item \textbf{Java JDK 17} y \textbf{Apache Maven} 3.9 o superior.
\item \textbf{Node.js} 18 LTS o superior y \textbf{Angular CLI} 21.
\item \textbf{Git} con configuración de cuenta en GitHub.
\item Conexión a Internet para la primera ejecución de \texttt{npm install} y
  \texttt{mvn} (descargan dependencias).
\end{itemize}

\section{Clonado del repositorio}

\begin{lstlisting}[language=bash, caption={Clonado del repositorio.}]
git clone https://github.com/amartii/TFG.git
cd TFG
\end{lstlisting}

\section{Datos de elevación}

El backend Java requiere un GeoTIFF configurado en \texttt{terrain.srtm-path}
que cubra el área de estudio. Su procedencia debe documentarse: un raster sintético
no es SRTM aunque un nombre o etiqueta lo indique. Los scripts MATLAB históricos
pueden utilizar un raster local; si no está disponible, algunas rutas recurren a
terreno alternativo o plano con advertencia. Debe comprobarse qué datos se utilizaron
efectivamente antes de interpretar una salida.

\section{Generación de simulaciones}

Se abre MATLAB en la carpeta \texttt{matlab/} y se ejecuta el \textit{script}
\texttt{generar\_escenarios}. El \textit{script} limpia la carpeta de salida y
genera tres ficheros JSON (URJC suburbano, Sierra rural y Madrid urbano) en
\texttt{matlab/output/}. Cada fichero contiene metadatos de la simulación y
hasta varios miles de puntos de cobertura.

\section{Backend Spring Boot}

Desde la raíz del proyecto:

\begin{lstlisting}[language=bash, caption={Arranque del backend.}]
cd backend
mvn spring-boot:run
\end{lstlisting}

Maven descarga las dependencias declaradas en \texttt{pom.xml} (Spring Boot~3.2.4,
\textit{springdoc-openapi}~2.3.0, Jackson, JUnit~5) y arranca el servidor
\textit{embedded} en el puerto~8080. Para verificar el correcto funcionamiento se
puede acceder a \texttt{http://localhost:8080/api/simulations} (lista de
simulaciones) o a \texttt{http://localhost:8080/swagger-ui.html} (documentación
interactiva).

\section{Frontend Angular}

En una segunda terminal:

\begin{lstlisting}[language=bash, caption={Arranque del frontend.}]
cd frontend
npm install     # solo la primera vez
npm start
\end{lstlisting}

El servidor de desarrollo de Angular escucha en el puerto~4200. Al abrir
\texttt{http://localhost:4200} en el navegador se carga la aplicación, que llama
al backend y muestra la primera simulación disponible.

\section{Docker Compose}
Existe una configuración de Docker Compose con raster sintético y persistencia
aislada. Se comprobó su configuración, pero la construcción y ejecución reales de
los contenedores y la verificación visual siguen pendientes.

\section{Preparación de cartografía vectorial (opcional)}

Se proporciona el \textit{script}
\texttt{scripts/download-pmtiles.ps1}, descrito en el capítulo~\ref{cap:capitulo4},
que permite preparar un fichero PMTiles local. El catálogo actual no ofrece temas
vectoriales, por lo que su descarga no activa el uso desconectado.

\chapter{Anexo C: Código Fuente Destacado}
\label{anx:codigo}

Este anexo recoge los fragmentos de código más representativos de cada capa del
sistema. La finalidad es proporcionar al lector un acceso rápido a la lógica
clave sin necesidad de navegar por el repositorio. El listado completo está
disponible en GitHub en el directorio raíz del TFG.

\section{Motor MATLAB --- Bloque 4 (modelo de propagación)}

\begin{lstlisting}[language=matlab, caption={Selección de modelo y cálculo de la pérdida con SRTM.}, label={cod:matlab_block4}]
% Auto-deteccion de Communications Toolbox
v_comm = ver('comm');
tiene_comm_toolbox = ~isempty(v_comm);

if tiene_comm_toolbox
    % Longley-Rice ITM via Communications Toolbox
    tx = txsite('Latitude', lat0, 'Longitude', lon0, ...
        'AntennaHeight', hTx_m, 'TransmitterFrequency', fc_Hz, ...
        'TransmitterPower', Ptx_W);
    pm = propagationModel('longley-rice');
    rx_sites = rxsite('Latitude', lats_flat, 'Longitude', lons_flat, ...
        'AntennaHeight', hRx_m);
    Prx_dBm = sigstrength(rx_sites, tx, pm);
    modelo_usado = 'Longley-Rice';
else
    % COST-231 Hata extendido + correccion de altura efectiva
    elev_tx = geointerp(elev_data, R, lat0, lon0, 'linear');
    hTx_eff = max(hTx_m + (elev_tx - elev_grid), 5);
    fc_MHz = fc_Hz / 1e6;
    a_hRx = (1.1*log10(fc_MHz) - 0.7)*hRx_m - (1.56*log10(fc_MHz) - 0.8);
    PL_dB = 46.3 + 33.9*log10(fc_MHz) - 13.82*log10(hTx_eff) - a_hRx ...
          + (44.9 - 6.55*log10(hTx_eff)) .* log10(dist_km) + 3;
    Prx_dBm = EIRP_dBm + Grx_dBi - PL_dB;
    modelo_usado = 'COST-231-Hata+SRTM';
end
\end{lstlisting}

\section{Backend Java --- validación de UUID y \textit{path traversal}}

\begin{lstlisting}[language=Java, caption={Lectura segura de un fichero JSON por UUID.}, label={cod:backend_service}]
private static final Pattern UUID_PATTERN = Pattern.compile(
    "^[0-9a-fA-F]{8}-[0-9a-fA-F]{4}-[0-9a-fA-F]{4}-" +
    "[0-9a-fA-F]{4}-[0-9a-fA-F]{12}$");

public JsonNode getById(String id) throws IOException {
    if (!UUID_PATTERN.matcher(id).matches()) {
        throw new SimulationNotFoundException(id);
    }
    Path file = dataDir.resolve(id + ".json").normalize();
    if (!file.startsWith(dataDir) || !Files.exists(file)) {
        throw new SimulationNotFoundException(id);
    }
    return cache.computeIfAbsent(id, k -> {
        try { return mapper.readTree(file.toFile()); }
        catch (IOException e) { throw new RuntimeException(e); }
    });
}
\end{lstlisting}

\section{Frontend Angular --- gestión de estado con \textit{signals}}

\begin{lstlisting}[language=JavaScript, caption={\texttt{SimulationService} con \textit{signals} y \texttt{computed}.}, label={cod:frontend_signals}]
private simulations = signal<Simulation[]>([]);
private selectedId = signal<string | null>(null);
private rsrpThreshold = signal<number>(-140);

readonly availableSimulations = this.simulations.asReadonly();
readonly currentThreshold = this.rsrpThreshold.asReadonly();

readonly currentSimulation = computed(() => {
  const id = this.selectedId();
  return this.simulations().find(s => s.simulation_id === id) ?? null;
});

readonly filteredCoveragePoints = computed(() => {
  const sim = this.currentSimulation();
  const threshold = this.rsrpThreshold();
  return sim?.coverage_points.filter(p => p.rsrp_dbm >= threshold) ?? [];
});
\end{lstlisting}

\section{Frontend Angular --- renderizado del \textit{heatmap}}

El listado~\ref{cod:heatmap} del capítulo~\ref{cap:capitulo4} muestra el cuerpo
de la función \texttt{updateHeatmap} que se ejecuta cada vez que cambian los
puntos filtrados.

\chapter{Anexo D: Especificación API REST}
\label{anx:api}

El backend Spring Boot expone una API REST documentada con OpenAPI~3.0 vía
\texttt{springdoc-openapi-ui} 2.3.0. La documentación interactiva está disponible
en tiempo de ejecución en \texttt{http://localhost:8080/swagger-ui.html} y la
especificación en formato JSON en \texttt{http://localhost:8080/api-docs}.

\section{Endpoints}

La API ofrece operaciones de consulta, exportación, creación y borrado bajo el
\textit{path} base \texttt{/api/simulations}:

\begin{table}[h!]
\centering
\caption{Endpoints expuestos por el backend.}
\label{tab:api_endpoints}
\small
\begin{tabular}{|l|l|p{6.5cm}|}
\hline
\textbf{Método} & \textbf{Ruta} & \textbf{Descripción} \\
\hline
GET & \texttt{/api/simulations}                & Lista resumida de simulaciones disponibles, opcionalmente paginada con \texttt{?page=0\&size=10}. \\
GET & \texttt{/api/simulations/\{id\}}         & JSON completo de una simulación, incluyendo todos los \texttt{coveragePoints}. \\
GET & \texttt{/api/simulations/\{id\}/export}  & Descarga el JSON con \texttt{Content-Disposition: attachment} para almacenamiento offline. \\
POST & \texttt{/api/simulations}              & Valida la petición, calcula una simulación con el motor Java y la persiste. \\
DELETE & \texttt{/api/simulations/\{id\}}      & Elimina la simulación por UUID; la interfaz solicita confirmación previamente. \\
\hline
\end{tabular}
\end{table}

\section{Códigos de respuesta}

\begin{itemize}
\item \texttt{200~OK}: respuesta correcta con cuerpo JSON.
\item \texttt{204~No Content}: borrado correcto, sin cuerpo de respuesta.
\item \texttt{400~Bad Request}: parámetros no válidos o UUID de borrado con formato incorrecto.
\item \texttt{409~Conflict}: el identificador de la ruta no coincide con la identidad del JSON.
\item \texttt{404~Not Found}: UUID no encontrado o fichero inexistente. La
  respuesta de error incluye \texttt{error} y \texttt{detail}.
\item \texttt{500~Internal Server Error}: errores inesperados de lectura, escritura o borrado.
\end{itemize}

\section{Esquema de la simulación}

El cuerpo de la respuesta de \texttt{GET /api/simulations/\{id\}} sigue el
esquema definido en el capítulo~\ref{cap:capitulo3} (sección de modelo de datos)
y comprende \texttt{id} (UUID v4), \texttt{name}, \texttt{createdAt}
(ISO~8601 UTC), \texttt{parameters} (parámetros de entrada), \texttt{metadata}
(metadatos descriptivos) y \texttt{coveragePoints} (lista de
\texttt{\{lat, lng, rsrp\}}). Los nombres de algunos parámetros difieren entre
escenarios MATLAB y Java.

\section{Ejemplos de uso con cURL}

\begin{lstlisting}[language=bash, caption={Listar simulaciones.}]
curl -X GET http://localhost:8080/api/simulations
\end{lstlisting}

\begin{lstlisting}[language=bash, caption={Obtener una simulación concreta (formato pretty con jq).}]
curl -s http://localhost:8080/api/simulations/a7ab9bad-3251-424e-915b-8ff3c532de9a \
     | jq '.metadata'
\end{lstlisting}

\begin{lstlisting}[language=bash, caption={Descargar el JSON como fichero.}]
curl -OJ http://localhost:8080/api/simulations/<UUID>/export
\end{lstlisting}

\begin{lstlisting}[language=bash, caption={Borrar una simulación tras confirmar la acción en la interfaz.}]
curl -X DELETE http://localhost:8080/api/simulations/<UUID>
\end{lstlisting}

\section{Política CORS}

El servicio configura una política CORS en la clase \texttt{CorsConfig} que permite
peticiones desde \texttt{http://localhost:4200} (servidor de desarrollo de Angular)
con los métodos GET, POST y DELETE. Spring gestiona las peticiones \textit{preflight}
OPTIONS para comprobar el método solicitado; OPTIONS no figura en la lista de métodos
de aplicación configurados. En producción este
\textit{origin} debe restringirse al dominio definitivo o eliminar el filtro si
backend y frontend se sirven desde el mismo \textit{host}.

\chapter{Anexo E: Resultados Adicionales}
\label{anx:resultados}
\textit{Escenarios rural montañoso y urbano denso.}

\chapter{Anexo F: Diagramas UML}
\label{anx:uml}

Este anexo recopila los diagramas UML que ilustran la estructura y el
comportamiento del sistema. Los diagramas están redactados con notación UML~2.5
y pueden compilarse a partir de las fuentes situadas en la carpeta
\texttt{TFG GD5G/figs/uml/} mediante la herramienta PlantUML.

\section{Diagrama de componentes}
\label{sec:uml_componentes}

El diagrama de componentes refleja la separación en tres capas que se describe
en los capítulos~\ref{cap:capitulo3} y~\ref{cap:capitulo4}: el motor MATLAB que
genera ficheros JSON, el backend Spring Boot que sirve esos ficheros mediante
una API REST y el frontend Angular que consume la API y representa la cobertura.
En la simulación bajo demanda el backend incorpora además el motor Java; MATLAB se
utiliza por separado para generar y contrastar referencias.

% \begin{figure}[h!]
% \centering
% \includegraphics[width=0.9\textwidth]{figs/uml/componentes.pdf}
% \caption{Diagrama de componentes del Gemelo Digital 5G.}
% \label{fig:uml_componentes}
% \end{figure}

\section{Diagrama de despliegue}
\label{sec:uml_despliegue}

El despliegue contempla tres nodos lógicos: el \textit{host} de simulación
(estación de trabajo con MATLAB), el servidor de aplicación (JVM ejecutando el
\textit{fat-jar} de Spring Boot) y el cliente (navegador con la SPA Angular). En
el escenario táctico previsto los tres nodos pueden coincidir físicamente en un
único portátil de campo, sin necesidad de red externa.

% \begin{figure}[h!]
% \centering
% \includegraphics[width=0.9\textwidth]{figs/uml/despliegue.pdf}
% \caption{Diagrama de despliegue mínimo (todo en un portátil).}
% \label{fig:uml_despliegue}
% \end{figure}

\section{Diagrama de clases del backend}
\label{sec:uml_clases}

Las clases principales del paquete \texttt{com.tfg.gemelo5g} son:
\texttt{SimulationController} (REST), \texttt{SimulationService} (lógica de
negocio y cache), \texttt{SimulationSummaryDto} (DTO de resumen),
\texttt{GlobalExceptionHandler} (gestión de errores) y
\texttt{SimulationNotFoundException}. La relación entre ellas es típica del
patrón \emph{Controller--Service}: el controlador inyecta el servicio y delega
la lógica de E/S y validación.

\section{Diagrama de secuencia: carga de una simulación}
\label{sec:uml_secuencia}

La secuencia de mensajes desde que el usuario selecciona una simulación hasta
que el \textit{heatmap} se renderiza es la siguiente:

\begin{enumerate}
\item Usuario selecciona simulación en el \texttt{<select>} del \textit{sidebar}.
\item \texttt{SidebarComponent} llama a
  \texttt{simulationService.selectSimulation(id)}.
\item \texttt{SimulationService} (Angular) emite \texttt{HttpClient.get('/api/simulations/\{id\}')}.
\item \texttt{SimulationController} (Spring) recibe la petición y delega en
  \texttt{SimulationService.getById(id)}.
\item \texttt{SimulationService} (Spring) valida el UUID, comprueba el
  \textit{path traversal}, consulta la cache y devuelve el \texttt{JsonNode}.
\item Spring serializa la respuesta como JSON con compresión GZIP.
\item Angular recibe el JSON, lo mapea a \texttt{Simulation} y actualiza el
  \textit{signal} \texttt{currentSimulation}.
\item El \texttt{effect()} de \texttt{MapComponent} se dispara, reconstruye el
  \textit{heatmap} con \texttt{leaflet.heat} y centra el mapa en la antena TX.
\end{enumerate}

% \begin{figure}[h!]
% \centering
% \includegraphics[width=\textwidth]{figs/uml/secuencia_carga.pdf}
% \caption{Secuencia de carga de una simulación.}
% \label{fig:uml_secuencia}
% \end{figure}
